\section{Глава~\ref{sec:code}. Азбука Морзе}

Предельные оценки главы --- теория информации и счёт; измерено, где в канале возникает шум, как быстро работает мозг и сколько стоит импульс, а вопрос о том, каким кодом кора пользуется на самом деле, остаётся открытым.

{\footnotesize
\setlength{\tabcolsep}{3pt}\renewcommand{\arraystretch}{1.15}
\begin{longtable}{>{\raggedright}p{0.5cm}>{\raggedright}p{4.0cm}>{\raggedright}p{5.6cm}>{\raggedright}p{2.3cm}>{\raggedright\arraybackslash}p{1.7cm}}
\toprule
№ & Утверждение & Опыт и что измерено & Источник & Статус \\
\midrule\endfirsthead
\toprule
№ & Утверждение & Опыт и что измерено & Источник & Статус \\
\midrule\endhead
1 & Частоты \nt{ГЛУТ}-нейронов коры --- от долей до десятков импульсов в секунду, и большую часть времени нейроны работают у нижней границы & Запись через прилегающую к клетке пипетку (выбор нейрона не зависит от его активности) в слуховой коре бодрствующей крысы: в каждый момент частоту выше $20$ импульсов в секунду дают меньше $5\,\%$ нейронов; у части нейронов фоновая частота ниже $0.01$ в секунду; распределение частот логнормальное & \cite{Hromadka2008} & измерено \\
2 & Ёмкость импульсного канала $R_{\max}\approx r\log_2\frac{e}{r\Delta t}$~\eqref{eq:spikeentropy}, почти вся она --- во времени импульсов (утверждение~\ref{prop:capacity}) & Энтропия двоичной строки из клеток длины $\Delta t$. Числа главы: $8$ бит на импульс при $r=10$ в секунду и $\Delta t=1\ms$, около $5$ бит при $\Delta t=10\ms$, меньше $2$ бит в секунду у счётчика импульсов & \cite{MacKay1952} & теория \\
3 & Чувствительные нейроны передают от одного до нескольких бит на импульс, в лучших случаях до половины предела & Восстановление раздражителя по импульсам чувствительных нейронов, для которых раздражитель известен; сводка измерений в книге & \cite{Rieke1997} & измерено \\
4 & Генератор импульсов точен; точное время задают быстрые перепады входа & Нейроны коры крысы в срезах: на повторяемый ток с быстрыми колебаниями импульсы повторяются с точностью меньше $1\ms$, на постоянный ток моменты расползаются & \cite{Mainen1995} & измерено \\
5 & Синапс ненадёжен: больше половины импульсов не доходят до получателя из-за случайности выброса & Минимальная стимуляция входов пирамидных нейронов гиппокампа (срезы): больше половины импульсов не вызывают ответа. Исключены колебания порога стимуляции, сбои проведения на ветвлениях аксона и низкая температура опыта; остаётся вероятностный выброс & \cite{AllenStevens1994} & измерено \\
6 & Поток импульсов в живой коре выглядит случайным: интервалы разбросаны как у пуассоновского потока, дисперсия числа импульсов близка к среднему; часть «шума» --- сообщение о движениях животного & Записи в зрительных областях V1 и MT бодрствующей обезьяны: коэффициент вариации интервалов около $1$, отношение дисперсии числа импульсов к среднему --- как у случайного процесса; модель, суммирующая независимые малые входы, даёт разброс гораздо меньше. В зрительной коре мыши значительная часть разброса предсказывается по видеозаписи движений самого животного & \cite{Softky1993,Stringer2019} & измерено \\
7 & Частотный код медленен: ошибка $1/\sqrt{rT}$~\eqref{eq:rateerror}, а мозг узнаёт картинку за $\sim150\ms$ & Формула --- пуассоновская статистика, вывод главы. Опыт: люди решали, есть ли животное на незнакомой фотографии, показанной на $20\ms$; потенциалы мозга на «да» и «нет» расходятся примерно через $150\ms$. Разбиение этого времени на десяток ступеней по $10$--$20\ms$ --- оценка главы, не измерение & \cite{Thorpe1996} & измерено \\
8 & Усреднение по группе ограничено корреляцией: ошибка падает не больше чем в $1/\sqrt c$ раз~\eqref{eq:correlated} (утверждение~\ref{prop:pooling}) & Пары нейронов области MT обезьяны, записанные одновременно: коэффициент корреляции чисел импульсов в среднем $0.12$, и теоретический анализ показывает, что даже такая корреляция ограничивает выигрыш от группы. Теория: предел ставит только та часть корреляции, что похожа на сигнал. Модель противоположной позиции: частота группы из $50$--$100$ нейронов читается за один межимпульсный интервал, $10$--$50\ms$, и время отдельных импульсов в коре не используется & \cite{Zohary1994,MorenoBote2014,ShadlenNewsome1998} & измерено \\
9 & Число можно писать во время первого импульса~\eqref{eq:latency}, в порядке импульсов группы или в фазе местного ритма & Сетчатка саламандры: пространственная структура коротко показанной картинки записана в относительных задержках первых импульсов выходных нейронов, почти независимо от контраста. Клетки места гиппокампа крысы: при проходе через «своё» место импульсы сдвигаются всё раньше в цикле тета-ритма ($7$--$12$\,Гц), сдвиг от $100^\circ$ до $355^\circ$. Насколько кора пользуется временем импульсов --- открытый вопрос (строка~8) & \cite{Gollisch2008,OKeefe1993} & измерено \\
10 & Какой код прочитан, решает постоянная времени получателя~\eqref{eq:reader} (утверждение~\ref{prop:reader}); в активной коре нейроны ближе к детекторам совпадений & Формула~\eqref{eq:reader} --- вывод главы. Обзор записей in vivo: в активной коре проводимость мембраны в несколько раз выше, чем в покое, и постоянная времени соответственно короче. Что кора переключает код именно так, прямо не показано & \cite{Destexhe2003} & гипотеза \\
11 & Истощающийся синапс передаёт изменение частоты, по существу $\Delta r/r$~\eqref{eq:depression} (рис.~\ref{fig:depression}) & Парные записи пирамидных нейронов коры: скорость истощения задаётся вероятностью выброса; при медленном истощении ответ отражает частоту, при быстром --- совпадения. Модель на основе измеренного истощения в коре крысы: одинаковые относительные изменения частоты на быстрых и медленных входах дают одинаковый ответ, то есть синапс передаёт $\Delta r/r$. Числа $1.7\to2.2$, $6.7$ и $90\ms$ --- расчёт главы & \cite{Tsodyks1997,Abbott1997depression} & измерено \\
12 & Пачка --- «тире»: её вероятность потеряться $(1-p)^k$~\eqref{eq:burst} мала, облегчение уменьшает её ещё & Обзор: на ненадёжных синапсах одиночные импульсы часто теряются, а пачки передаются надёжно благодаря облегчению; пачка вызывает долговременные изменения синапсов. Что мозг читает пачку как отдельный знак, --- гипотеза & \cite{Lisman1997,AllenStevens1994} & гипотеза \\
13 & Быстрые \nt{ГАМК}-нейроны задают ритм и «пунктуацию»; ещё один класс тормозит тормозящих и открывает ворота (утверждение~\ref{prop:punctuation}) & Обзор свойств классов \nt{ГАМК}-нейронов коры. Оптогенетика: ритмическое возбуждение быстрых \nt{ГАМК}-нейронов вызывает гамма-ритм, и ответ на раздражитель зависит от фазы цикла. В зрительной коре мыши PV-нейроны тормозят друг друга, SST-нейроны --- все остальные классы, VIP-нейроны --- преимущественно SST. Возбуждение VIP-нейронов у бодрствующей мыши снимает торможение с \nt{ГЛУТ}-нейронов; их включает сигнал подкрепления. Разделение ролей как общее правило --- гипотеза & \cite{Tremblay2016,Cardin2009,Pfeffer2013,Pi2013} & гипотеза \\
14 & Энергия ограничивает среднюю частоту величиной порядка одного импульса в секунду~\eqref{eq:energy} & Энергетический бюджет серого вещества грызуна по анатомическим и физиологическим данным: импульсы и синаптические токи --- $47\,\%$ и $34\,\%$ энергии сигнализации; одновременно могут быть активны не больше $\sim15\,\%$ нейронов. Для коры человека: не больше $\sim1\,\%$ нейронов могут быть сильно активны одновременно. Число $\sim10^9$ АТФ на импульс и мощность $\sim10$\,Вт на сигнализацию --- оценка главы на основе этих расчётов & \cite{Attwell2001,Lennie2003} & теория \\
15 & Код разрежен и подстроен под наибольшее число бит на джоуль; кора меняет точность на энергию (утверждение~\ref{prop:economy}) & Гипотеза экономного кода. Опора: у мышей при ограничении пищи в зрительной коре падает проводимость AMPA-рецепторов и синаптический расход АТФ ($-29\,\%$), частота импульсов сохраняется, а настройка на ориентацию расширяется на $32\,\%$ & \cite{Barlow1961,Padamsey2022} & гипотеза \\
\bottomrule
\end{longtable}}
